\abstract{Large language models can emit fluent text while leaving intermediate semantic structure implicit. We study whether explicit event-role and logical-primitive side-state can improve a pretrained decoder without damaging its language behavior. We introduce PAT-ER, a decoder architecture with a normal token stream, an event-role register stream, and a primitive register stream. The primitive stream is motivated by the view that logical primitives answer characteristic semantic questions, such as what licenses a conclusion, what conflicts with it, or why evidence is insufficient. Across eight seeds on the same Qwen3-0.6B backbone, replacing token-pooled auxiliary heads with typed PAT-ER registers improves primitive macro-F1 by 0.209 (95\% CI [0.182, 0.237]) and role-to-primitive macro-F1 by 0.091 (95\% CI [0.074, 0.110]) with no language-model loss cost. A generic-register control shows that this is not merely the effect of adding latent registers: typed PAT-ER improves over generic registers by 0.116 primitive macro-F1 and 0.110 role-to-primitive macro-F1, with both confidence intervals excluding zero. A warm-started model then recovers pretrained language quality (LM loss 1.344 versus 2.555 for the frozen-backbone register model) while retaining most side-state behavior. Finally, a decoupled interface mode produces robust schema-grounded function calls on 242 held-out prompts (Hermes parse 0.952, exact arguments 0.981, JSON validity 1.000, IDK F1 1.000) while base-mode side-state metrics remain byte-identical to the warm-start baseline. The model is not a theorem prover and does not achieve perfect unseen tool-name copying; the contribution is a measured architecture signal and a usable, guarded interface.}

\keyword{large language models; neural-symbolic learning; semantic roles; logical reasoning; tool use; model architecture; function calling}
